\documentclass{article}
\usepackage[utf8]{inputenc}
\usepackage[T1]{fontenc}
\usepackage{graphicx}
\usepackage{amsmath}
\usepackage{amssymb}
\usepackage{booktabs}
\usepackage{multirow}
\usepackage{float}
\usepackage{tabularx}
\usepackage[section]{placeins}
\usepackage{pdflscape}
\usepackage{longtable}
\usepackage{xurl}
\usepackage[
  colorlinks=true,
  linkcolor=blue,
  urlcolor=blue,
  citecolor=blue,
  bookmarks=true,
  bookmarksopen=true,
  pdftitle={Early Detection of Illegal Territorial Activities via Satellite Imagery},
  pdfauthor={Vélez Muñoz, Gonzalez Veloza, Ahumada León}
]{hyperref}
\usepackage{geometry}
\geometry{margin=1in}

% ----------------------------------------------------------------
\title{%
  Early Detection of Illegal Territorial Activities\\
  via Satellite Imagery: A Cloud-Based\\
  Multi-Phenomenon Monitoring Architecture
}
\author{
  Juan Pablo Vélez Muñoz \\
  Luiza Mariana Gonzalez Veloza \\
  Daniel Alexander Ahumada León
}
\date{November 2026}
% ----------------------------------------------------------------

\begin{document}

\begin{titlepage}
    \centering
    \vspace*{2cm}
    {\Large \textbf{Early Detection of Illegal Territorial Activities\\
    via Satellite Imagery: A Cloud-Based\\
    Multi-Phenomenon Monitoring Architecture}\par}
    \vspace{1.5cm}
    {\Large Enterprise Architecture and Digital Transformation \par}
    \vspace{2cm}
    {\large
    Juan Pablo Vélez Muñoz \\
    Luiza Mariana Gonzalez Veloza \\
    Daniel Alexander Ahumada León
    \par}
    \vfill
    {\large November 2026 \par}
\end{titlepage}

\tableofcontents
\newpage

% =============================================================
% ABSTRACT
% =============================================================
\begin{abstract}
In Colombia, deforestation, illegal gold mining, and coca cultivation are monitored
through separate institutional pipelines that operate on quarterly or annual cycles,
creating a gap of months to a year between when territorial changes occur and when
authorities can respond.
Satellite imagery provides the raw data necessary to close this gap, but no existing
system jointly detects these three phenomena with sufficient accuracy and operational
frequency to support near-real-time institutional response.
The closest precedent---Riocampo et al.\ (2023)---achieved only 60--70\% detection
rates using optical imagery alone, without addressing false-alarm management or
cloud-cover robustness.
This project proposes a cloud-based, scalable monitoring architecture that ingests
multi-source satellite data (optical, with SAR radar as an optional module), applies a multi-label
classification pipeline to jointly detect the three phenomena, and exposes alerts
through a structured API consumed by institutional stakeholders.
The architecture is designed for deployment at national scale, with components
spanning automated satellite ingestion, preprocessing, multi-source feature fusion,
joint classification, false-alarm filtering, and alert delivery to existing
government information systems.
The central design hypothesis is that jointly detecting the three phenomena with
a shared multi-label model, exploiting their geographic co-occurrence, lowers
false alarms relative to three independent detectors.
Testing it requires spatial ground truth from IDEAM and UNODC/SIMCI, which we
could not obtain within the project timeline.
We therefore adjusted the research question and the minimum viable product to
openly available data: a pilot on forest-loss detection in Caquetá and Putumayo
(Hansen/GLAD labels) that asks whether adding Sentinel-1 radar to Sentinel-2
imagery improves detection and robustness to cloud occlusion.
Optical plus SAR input matched optical-only detection on clean imagery (mean
average precision 0.778 vs.\ 0.776 over three seeds) and showed a small
advantage under synthetic occlusion (+0.022 average precision, in all three
seeds), but recall (0.64) fell short of the 0.75 criterion, so the adjusted
question was not answered in favor of fusion.
The original hypothesis remains viable and untested: since the pilot could not
demonstrate a benefit from radar, we defer it on cost grounds and propose that
the first test use optical imagery only, keeping SAR as an optional module.
The FastAPI inference service was deployed to Google Cloud Run on 17 November 2026; two reference requests reproduced the Colab probabilities, without constituting a new model evaluation.
The contribution is a reproducible reference architecture, a working
detection pipeline, and an honest account of what the available data could and
could not establish.
\end{abstract}

\newpage

\section*{Executive Summary}
\addcontentsline{toc}{section}{Executive Summary}

\paragraph{Problem.}
Colombian authorities learn of deforestation, illegal alluvial mining, and coca
cultivation months to a year after they happen, because each phenomenon is
monitored through a separate institutional pipeline.

\paragraph{Proposal.}
Centinel is a cloud architecture that ingests satellite imagery, detects
territorial change, filters false alarms, and delivers georeferenced alerts
through an API. Its central design hypothesis is that detecting the three
phenomena jointly, with a shared multi-label model, lowers false alarms
compared with three independent detectors.

\paragraph{What we did.}
The official ground truth needed to test that hypothesis (IDEAM, UNODC/SIMCI,
EVOA) could not be obtained in time. We built a working pipeline on open data
and used it to ask a related question: does adding Sentinel-1 radar to
Sentinel-2 imagery improve annual forest-loss detection? The pilot covers
about 18{,}000\,km$^2$ of Caquet\'a and Putumayo and 11{,}764 image patches. The FastAPI container was deployed on Google Cloud Run and checked on 17 November 2026: two reference queries reproduced the Colab outputs (0.9993 and 0.0338). This is a deployment check, not a new model evaluation.

\paragraph{What we found.}
Radar gave no meaningful gain on clean imagery (average precision 0.778 vs.\
0.776) and a small gain under synthetic cloud occlusion (+0.022). Recall was
0.64 against a target of 0.75, so the pilot did not meet its own criteria.
The central hypothesis was not tested and remains viable.

\paragraph{Cost.}
Measured and derived estimates indicate that compute and storage are small
(order of tens to hundreds of US dollars per year for one national annual
pass). The costs that matter are the Earth Engine license status
(about 6{,}000 US dollars per year if the service is classed as commercial),
the services we did not price (database, orchestration, observability), and
staff. The Cloud Run demonstration tests were covered by Google Cloud trial credit, with a USD~1 budget alert; no monthly operating cost was measured.

\paragraph{Main risks.}
Two have already occurred: missing ground truth and recall below target.
The remaining ones concern whether the labels reflect illegality, licensing,
false-alarm trust, generalization to other regions, and the governance of
sensitive location data. The demonstration endpoint is public and unauthenticated; production access controls are not implemented.

\paragraph{Recommendation.}
Neither scale nor stop. Proceed to a conditional second phase whose first step
is to obtain official ground truth and test the joint-versus-independent
hypothesis with optical imagery only. Do not commit to radar ingestion or to
national coverage until that test and a regional deployment succeed.

\newpage


% =============================================================
% INTRODUCTION
% =============================================================
\section{Introduction}

\subsection{Context}

Monitoring large remote territories requires substantial institutional resources
that Colombia's environmental and judicial agencies do not consistently have
available.
Deforestation, illegal gold mining (EVOA), and coca cultivation are currently
tracked through separate pipelines: IDEAM's SMBYC publishes deforestation alerts
quarterly; UNODC and SIMCI release coca cultivation figures annually; and EVOA
reporting operates on a similarly long cycle.
Indigenous communities and \textit{resguardos} in affected regions---often the
first to perceive on-the-ground changes---and national bodies such as IDEAM, the
Ministry of Environment, and the security forces may not receive official alerts
until several months or up to a year after the damage occurs~\cite{mongabay_siona2023}.

Critically, these three phenomena do not occur in isolation.
UNODC data show that alluvial gold mining and coca cultivation frequently overlap
in the same territories~\cite{unodc_evoa2022}; their relation to deforestation is
more nuanced (Section~\ref{sec:problem}), which points to a partial co-occurrence
that the current siloed monitoring approach does not exploit.

\subsection{Problem}

Satellite imagery is already freely and continuously available, so the core
bottleneck is not data access but the capacity to process that data at operational
frequency and at the required spatial precision.
The persistent monitoring gap---between the moment a territorial change occurs and
the moment an authority can act on it---creates a window in which evidence
disappears, perpetrators relocate, and ecological damage compounds.

This raises a concrete technical question: \textit{is it better to continue
monitoring each phenomenon separately, given their distinct spectral and
structural signatures in satellite imagery, or is it feasible and advantageous to
detect them jointly, leveraging their geographic co-occurrence?}
Any viable answer must also manage false alarm rates carefully, because
institutional trust in an automated alert system depends on its specificity,
not only its sensitivity.

\subsection{Contributions}

This work makes the following contributions at this stage of the project:

\begin{itemize}
    \item Using official data from IDEAM and UNODC, we document that illicit
    mining and coca cultivation overlap in a large share of the territories where
    either occurs, that their overlap with deforestation is partial, and that the
    areas of coca and illicit mining have grown while national deforestation
    rebounded in 2024 and 2025 after its 2023 minimum.
    \item We conduct a critical comparative review of existing satellite monitoring
    systems---SMBYC, Global Forest Watch (GLAD and RADD alerts), CoMiMo, and Riocampo et al.\ (2023)---and demonstrate that none of them simultaneously addresses
    multi-phenomenon joint detection, cloud-cover robustness, and false-alarm
    management at operational scale.
    \item We propose a large-scale cloud-based monitoring architecture whose
    principal components include an automated satellite ingestion pipeline,
    multi-source feature fusion (optical, with SAR as an optional module), a joint multi-label
    classification layer, a false-alarm filtering stage, and a structured
    alert delivery API consumable by existing government platforms.
    \item We report a reduced proof-of-concept pilot. Because official
    ground truth for the central hypothesis (joint versus independent detection)
    could not be obtained in time, the pilot tests a related, data-feasible
    question---whether SAR improves forest-loss detection---with explicit
    go/no-go criteria. The criteria were not met, and we derive from this how
    the original hypothesis should be tested next.
\end{itemize}

\subsection{Document Structure}

The remainder of this document is organized as follows.
Section~\ref{sec:problem} characterizes the problem with qualitative and
quantitative evidence.
Section~\ref{sec:sota} reviews existing systems and their limitations.
Section~\ref{sec:gap} identifies the residual gap and presents our proposed
solution and architecture direction.
Section~\ref{sec:architecture} describes the proposed system architecture.
Section~\ref{sec:prototype} presents the prototype implementation.
Section~\ref{sec:results} reports the results obtained.
Section~\ref{sec:business} addresses cost, governance, security, risk, and a
roadmap.
Finally, Section~\ref{sec:conclusions} summarizes conclusions and outlines
future work.


% =============================================================
% PROBLEM
% =============================================================
\section{Problem}
\label{sec:problem}

\subsection{Qualitative Aspects}

Indigenous communities in the Colombian Amazon are typically the first to perceive
the effects of deforestation and illicit crops, often before any official alert
is issued.
Young members of the Siona people, for example, have publicly reported the loss
of fauna and the disruption of their spiritual relationship with their territory
due to encroaching deforestation and illegal activity~\cite{mongabay_siona2023}.

State interventions, when they occur, do not consistently produce lasting results.
Conservation specialists have documented that a law enforcement operation against
illegal mining in Parque Nacional Río Puré had measurable effects for only
approximately two months, after which activity resumed due to the absence of
sustained institutional presence in remote areas~\cite{mongabay_riopure2026}.
This pattern illustrates a systemic challenge: without continuous, automated
monitoring, illegal activity relocates rather than disappears, and each response
cycle begins from scratch.

\subsection{Quantitative Aspects}

Official figures show rising areas of coca cultivation and illicit mining and a rebound in deforestation after its 2023 minimum, summarized
in Table~\ref{tab:quant}.

\begin{table}[ht]
\centering
\caption{Trends in deforestation, coca cultivation, and illegal gold mining in Colombia}
\label{tab:quant}
\resizebox{\textwidth}{!}{%
\begin{tabular}{llrrl}
\toprule
\textbf{Phenomenon} & \textbf{Year} & \textbf{Area (ha)} & \textbf{Change} & \textbf{Source} \\
\midrule
\multirow{3}{*}{Deforestation}
 & 2023 & 79{,}256  & ---    & IDEAM / SMByC~\cite{ideam_smbyc} \\
 & 2024 & 113{,}608 & +43\%  & IDEAM / SMByC \\
 & 2025 & 119{,}483 & +5.2\% & El Tiempo~\cite{eltiempo_deforestacion2025} \\
\midrule
\multirow{3}{*}{Coca cultivation}
 & 2022 & 230{,}000 & ---    & UNODC / SIMCI \\
 & 2023 & 253{,}000 & +10\%  & UNODC / SIMCI~\cite{minjusticia_coca2023} \\
 & 2024 & 261{,}000 & +3.5\% & UNODC / SIMCI~\cite{simci2024} \\
\midrule
\multirow{3}{*}{Alluvial gold mining (EVOA)}
 & 2023 & 105{,}060 & ---    & UNODC / MinEnergía~\cite{minenergia_evoa2025} \\
 & 2024 & 113{,}310 & +7.9\% & UNODC / MinEnergía \\
 & 2025 & 118{,}410 & +4.5\% & UNODC / MinEnergía~\cite{minenergia_evoa2025} \\
\bottomrule
\end{tabular}%
}
\vspace{2mm}
\begin{flushleft}
\footnotesize
\textit{Note:} In 2023, deforestation reached its lowest level in 23 years before
rising sharply in 2024.
In 2025, 75\% of the area with evidence of alluvial gold exploitation
(118{,}410\,ha) was classified as illicit, about 88{,}860\,ha~\cite{minenergia_evoa2025}.
The UNODC coca report for 2025 had not been published at the time of writing, so
the UNODC series ends in 2024 (the SIMCI 2024 report appeared in June 2026).
The National Police's SIIMA system, which uses its own methodology, reports
roughly 258{,}000--262{,}000\,ha for 2025 depending on the cut-off date
(August or December)~\cite{siima2025}; the two series are not directly comparable.
\end{flushleft}
\end{table}

The relationship among these phenomena is not uniform but is structurally
significant.
The UNODC EVOA 2025 report found that about 45\% of the area with evidence of
alluvial gold exploitation on land coincided with areas that had coca cultivation
in 2024~\cite{minenergia_evoa2025}; the EVOA 2022 report had already shown
overlap in 76 of 101 mining municipalities~\cite{unodc_evoa2022}.
However, the relationship with deforestation is more nuanced: Colombia's Ministry of Environment estimates
(2023) that only approximately 13\% of annual deforestation is linked to illicit crops~\cite{minambiente_13},
while IDEAM reports identify pasture conversion for land grabbing and extensive
cattle ranching among its primary direct causes.
This heterogeneity is precisely what makes joint detection both challenging and
informative: a system that can distinguish causes of land-cover change, rather
than merely detecting that change occurred, provides substantially more actionable
intelligence to responding institutions.


% =============================================================
% STATE OF THE ART
% =============================================================
\section{State of the Art}
\label{sec:sota}

\subsection{Grey Literature}

The \textbf{SMByC} (Sistema de Monitoreo de Bosques y Carbono) of IDEAM has
published quarterly Early Deforestation Alerts (ATD) using MODIS and Landsat
imagery since 2016~\cite{ideam_smbyc,smbyc_atd_github}.
The \textbf{UNODC}, through \textbf{SIMCI}, produces annual coca cultivation
monitoring reports~\cite{minjusticia_coca2023}, and through its \textbf{EVOA}
reports tracks illegal alluvial gold extraction~\cite{unodc_evoa2022}.
Investigative journalism outlets such as \textbf{Mongabay Latam} complement these
institutional sources by documenting community-level effects and the practical
limitations of state intervention in remote
regions~\cite{mongabay_siona2023,mongabay_riopure2026}.

\subsection{Scientific Literature}

\textbf{Saavedra-Romero (2017)} developed an early change-detection model for
forest cover achieving 79\% accuracy, which served as the methodological
foundation for CoMiMo~\cite{comimo_urosario}---a system for detecting illegal
mining developed at Universidad del Rosario.
CoMiMo demonstrated the feasibility of near-real-time satellite-based detection
for a single phenomenon, but its scope is restricted to mining and it has not
been extended to the multi-phenomenon case.

The most directly related work is \textbf{Riocampo et al.\ (2023)}~\cite{riocampo2023},
who attempted to jointly detect illegal mining, deforestation, and illicit crops
from satellite imagery using convolutional neural networks applied to optical
imagery.
Their system achieved a detection rate of 60--70\%, which falls below the
operational threshold that would be needed for institutional deployment.
Critically, the study does not report precision, false positive rates, or
any mechanism for managing false alarms---an omission that is especially
significant given that false alerts erode institutional trust in automated
systems and can trigger costly on-the-ground verifications.
Additionally, Riocampo et al.\ relied exclusively on optical imagery, which is
vulnerable to cloud cover: in Colombia's tropical regions, cloud obstruction
can persist for weeks at a time, creating blind spots precisely during
periods of active environmental change.
Their work also does not address the question of scalability or operational
deployment, making it a proof-of-concept that does not translate directly
to an actionable monitoring platform.

The \textbf{FACSAT-1} program of the Colombian Air Force demonstrated the use
of neural networks for open-pit mining pattern detection~\cite{facsat1}, but
this work targets a narrow sub-problem and does not generalize to the
multi-phenomenon, multi-source architecture that an operational system
would require.

\subsection{Technical Documentation}

The SMByC's ATD pipeline~\cite{smbyc_atd_github} uses MODIS for high-frequency,
low-resolution coverage and Landsat for high-resolution but less frequent updates.
This dual-sensor design illustrates the classic trade-off between temporal
resolution and spatial resolution in satellite monitoring, and it motivates the
addition of SAR (Synthetic Aperture Radar) imagery---which is available at
competitive temporal resolution and is not affected by cloud cover---as a third
data source in our proposed architecture.

SAR is not new in operational deforestation monitoring.
The RADD alerts, produced by Wageningen University from Sentinel-1 and
available on Global Forest Watch, detect forest disturbances in the humid
tropics in near real time and have been extended to the Amazon
basin~\cite{reiche2021radd,gfw_radd}; Global Forest Watch's integrated alerts
combine optical and radar systems~\cite{gfw_integrated}.
Research on optical--SAR fusion for deforestation under cloud cover in the
Brazilian Amazon reports that fusion improves detection~\cite{ferrari2023}.
These systems and studies address forest disturbance only; what remains open,
and what this work targets, is the joint detection of deforestation, illegal
mining and coca cultivation together with false-alarm management and delivery
to institutional consumers.

\subsection{Comparative Summary}

Table~\ref{tab:sota} summarizes the capabilities and limitations of the most
relevant existing systems.

\begin{table}[ht]
\centering
\caption{Comparison of existing satellite monitoring systems}
\label{tab:sota}
\resizebox{\textwidth}{!}{%
\begin{tabular}{lccccl}
\toprule
\textbf{System} & \textbf{Deforestation} & \textbf{Mining} & \textbf{Coca} & \textbf{SAR} & \textbf{Key Limitation} \\
\midrule
SMByC / ATD         & \checkmark & --         & --         & -- & Quarterly cadence; causes described in reports, not per detection \\
GLAD (GFW)          & \checkmark & --         & --         & -- & Optical alerts; broad driver classes, not legality; no mining or coca \\
RADD (GFW)          & \checkmark & --         & --         & \checkmark & Radar-only alerts; detects disturbance, not legality; no mining or coca \\
CoMiMo              & --         & \checkmark & --         & -- & Single phenomenon; not publicly deployed \\
FACSAT-1            & --         & \checkmark & --         & -- & Narrow scope; no multi-label output \\
Riocampo et al.\ (2023) & \checkmark & \checkmark & \checkmark & -- & 60--70\% recall; no false-alarm management; optical only \\
\textbf{This work (proposed)} & \checkmark & \checkmark & \checkmark & (\checkmark) & Proposed architecture; SAR optional; pilot covers deforestation only \\
\bottomrule
\end{tabular}%
}
\end{table}


% =============================================================
% MOTIVATION / GAP / SOLUTION
% =============================================================
\section{Motivation, Gap, and Proposed Solution}
\label{sec:gap}

\subsection{Identified Gap}

The comparative review in Section~\ref{sec:sota} reveals a gap with two
interlocking dimensions.
The first is \textit{coverage}: no currently operational system jointly detects
all three phenomena---deforestation, illegal mining, and coca cultivation---with
sufficient accuracy for institutional use.
The second is \textit{operationalizability}: even the only prior work that
attempted joint detection (Riocampo et al., 2023) is a standalone
proof-of-concept that does not address cloud-cover robustness, false-alarm
management, alert delivery, or deployment at scale.

\subsection{Proposed Solution}

We propose a \textbf{cloud-based, multi-source satellite monitoring architecture}
designed to continuously ingest satellite data, jointly classify territorial
changes by phenomenon type, filter false alarms, and deliver structured alerts
to institutional consumers via an API.

The principal design decisions that distinguish this proposal from prior work are:

\begin{enumerate}
    \item \textbf{Multi-source fusion:} The architecture would ingest both
    optical imagery (Sentinel-2, Landsat) and SAR imagery (Sentinel-1) to
    maintain observability during cloud-cover periods, which Riocampo et al.\
    (2023) do not address. Radar alone is not new in operational monitoring
    (RADD alerts~\cite{reiche2021radd}); what this proposal adds is its
    combination with optical data inside a joint detector of three phenomena. The pilot (Section~\ref{sec:results}) did not show a
    decisive benefit of SAR on annual composites, so in the first validation of
    the joint hypothesis SAR is treated as an optional module.
    \item \textbf{Multi-label classification:} Rather than operating three
    independent binary classifiers in isolation, the proposed system would
    evaluate whether a shared multi-label model trained on the joint
    distribution of the three phenomena achieves better calibration and lower
    false alarm rates, particularly in zones of co-occurrence. This is the
    central architectural hypothesis of the proposal. The pilot reported here
    does not test it, because the required ground truth could not be obtained
    in time; Section~\ref{sec:conclusions} describes how it should be tested
    next.
    \item \textbf{Explicit false-alarm management:} A post-classification
    filtering stage would apply temporal consistency checks and spatial
    context analysis to suppress spurious alerts, with a tunable
    precision-recall trade-off exposed to institutional operators through
    configurable thresholds.
    \item \textbf{Structured alert delivery API:} Rather than producing image
    outputs requiring manual inspection, the system would expose a REST API
    delivering structured, georeferenced alerts consumable by existing
    institutional platforms such as SMBYC dashboards and government data
    portals.
    \item \textbf{Cloud-native scalability:} The architecture would be built
    on cloud-native principles---containerized microservices, managed object
    storage for satellite data, auto-scaling inference endpoints---so that
    geographic coverage could be expanded from a pilot region to the national
    territory without proportional increases in operational cost or manual
    intervention.
\end{enumerate}

\subsection{Proof-of-Concept Pilot and Decision Criteria}
\label{sec:pilot-criteria}

The full architecture is intended for large-scale deployment with the
institutional and computational resources of an agency such as IDEAM or a
national environmental technology fund.
As a first validation step, the proposal included a proof-of-concept pilot.

\paragraph{Original plan and data constraint.}
The pilot was originally designed to test the central hypothesis directly:
whether a joint multi-label classifier outperforms three independent detectors
in a region where deforestation, illegal alluvial mining, and coca cultivation
co-occur.
That test needs spatial ground-truth layers from IDEAM (official deforestation),
UNODC/SIMCI (coca cultivation), and the national alluvial-mining evidence
(EVOA).
We could not obtain these layers within the project timeline, and without them
the hypothesis cannot be evaluated.

\paragraph{Adjusted question.}
Rather than present an unevaluated design, we adjusted the question and the
minimum viable product to the data that were openly available: Sentinel imagery
and the global Hansen / GLAD forest-loss product.
The adjusted pilot asks a related, verifiable question about the architecture's
data-fusion decision: whether adding Sentinel-1 radar to Sentinel-2 optical
imagery improves the detection of annual forest loss, and whether it makes
detection more robust when part of the optical image is obscured, as is common
in Colombia's tropical regions.
It is therefore restricted to deforestation and does not test joint detection.

The adjusted pilot is structured as follows:
\begin{itemize}
    \item \textbf{Study area:} A region within Caquetá or Putumayo, Colombian
    departments with documented forest loss.
    \item \textbf{Data:} Annual Sentinel-2 (optical) and Sentinel-1 (SAR)
    composites for 2021--2023, paired as (previous year, evaluated year), with
    global forest-loss labels from the Hansen / GLAD dataset~\cite{hansen2013}
    as ground truth.
    \item \textbf{Experimental comparison:} Two configurations of the same
    network trained on identical data splits---(a) optical imagery only and
    (b) optical plus SAR imagery.
    \item \textbf{Metrics:} Precision, recall, F1 score, false positive rate,
    and average precision, reported both on the standard test set and under
    synthetic cloud occlusion of the optical channels.
\end{itemize}

The fusion approach was to be considered worth adopting within the full
platform if: (a) it achieves recall $\geq 0.75$, precision $\geq 0.70$, and
false positive rate $\leq 0.20$ on the test set; and (b) it improves average
precision over the optical-only configuration, on average and in at least two
of three random seeds.
Its behavior under
occlusion is a secondary outcome.
The thresholds in (a) are design targets that we set without an empirical basis
from prior work, and criterion (a) measures the absolute performance of the
fusion model, whereas only criterion (b) isolates the contribution of radar;
a better-designed test would have kept the two questions separate.
Section~\ref{sec:results} reports that these criteria were not met, and
Section~\ref{sec:conclusions} draws the consequence for the original hypothesis.


% =============================================================
% ARCHITECTURE
% =============================================================
\section{Architecture}
\label{sec:architecture}

The proposed system is described using the C4 model at three levels of
abstraction: the system context (Level~1), the container decomposition
(Level~2), and the component breakdown of the detection core (Level~3), which
contains the two containers that carry the contribution of this proposal: the
ML inference service and the false-alarm filter.
In all diagrams, an arrow points in the direction of the call or data flow, and
its label names the event or operation; blue boxes are parts of Centinel,
gray boxes are external systems, and dark blue boxes are people.

% ----------------------------------------------------------------
\subsection{Level 1 --- System Context}

Figure~\ref{fig:c4-l1} shows Centinel as a single box between the actors that
supply it with data and those that consume its output.
On the supply side, it pulls Sentinel-1 and Sentinel-2 scenes from the
\textbf{Copernicus Data Space}, Landsat scenes from \textbf{USGS / NASA},
reference forest-loss labels from \textbf{UMD / Global Forest Watch}, and
official ground-truth layers from \textbf{IDEAM--SMByC} and
\textbf{UNODC / SIMCI} (coca and alluvial mining) for training and validation.
On the consumption side, five types of people interact with the platform:
environmental analysts (IDEAM, MinAmbiente) and illicit-economies analysts
(UNODC / SIMCI, Ministry of Justice) review and query alerts; enforcement
officers (security forces, Fiscal\'ia, regional environmental authorities)
receive prioritized alerts; community monitors of indigenous territories report
what they find on the ground; and platform operators run the pipelines.
Alerts also leave the platform through a government interoperability platform
and a notification gateway (SMS, WhatsApp, e-mail).

\begin{figure}[H]
\centering
\includegraphics[width=\textwidth]{Cetinel_TDSE-C4_L1_Context.png}
\caption{C4 Level 1: system context of the Centinel platform.}
\label{fig:c4-l1}
\end{figure}

This boundary makes the system's responsibilities explicit: it does not produce
new satellite data, nor does it make enforcement decisions.
Its sole function is to reduce the latency between a territorial change occurring
and an institutional consumer receiving an actionable, reliable alert about it.

% ----------------------------------------------------------------
\subsection{Level 2 --- Container Decomposition}

The platform is decomposed into seventeen containers, shown in
Figure~\ref{fig:c4-l2} and described in Table~\ref{tab:containers}.
To keep the diagram readable, the people already shown in Level~1 do not
reappear; only the external systems that touch a container are drawn.
The figure reads in three zones.

\paragraph{Detection pipeline (top rows).}
The pipeline follows one satellite scene from arrival to alert.
A \textbf{workflow orchestrator} schedules the \textbf{scene discovery} service,
which searches the external satellite catalogs and publishes a
\textit{SceneFound} event; the \textbf{ingestion service} downloads the scene,
stores it in the \textbf{imagery data lake} and publishes \textit{SceneStored}.
The \textbf{preprocessing service} turns the raw scene into analysis-ready tiles
(cloud masking, calibration, co-registration), writes them to the lake and
publishes \textit{TileReady}.
The \textbf{ML inference service} consumes that event and emits detections, which
the \textbf{false-alarm filter} validates before writing confirmed alerts to the
\textbf{alert database}; each confirmed alert triggers the \textbf{notification
service}, which delivers it through the external notification gateway.
Because the stages communicate through events instead of direct calls, each one
can scale, fail and be retried independently.

\paragraph{MLOps loop (left, middle rows).}
The \textbf{labeling service} harmonizes the official ground-truth layers into
versioned label sets; the \textbf{training pipeline} combines them with tiles from
the data lake, trains and evaluates models, and registers them in the
\textbf{model registry}, from which the inference service loads the production
model.
The \textbf{monitoring service} promotes a model or triggers retraining based on
the outcomes verified in the field, which it obtains through the alert API.
This loop lets field verdicts improve future models; it is not covered by the
pilot of Section~\ref{sec:prototype}.

\paragraph{Access layer (right, lower rows).}
Analysts use the \textbf{web dashboard} and community monitors use an
offline-first \textbf{field verification app}; both, together with the
government interoperability platform, reach the platform only through the
\textbf{API gateway}.
The gateway validates every token against the government identity provider,
applies per-institution roles, and routes authorized calls to the \textbf{alert
API}, the only container that reads and writes the alert database on behalf of
users.

\begin{landscape}
\thispagestyle{empty}
\begin{figure}[H]
\centering
\includegraphics[width=\linewidth,height=0.78\textheight,keepaspectratio]{Cetinel_TDSE-C4_L2_Containers.png}
\caption{C4 Level 2: container decomposition of the Centinel platform.}
\label{fig:c4-l2}
\end{figure}
\vfill
\centering\thepage
\end{landscape}

{\small
\begin{longtable}{@{}p{3.3cm}p{3.7cm}p{8cm}@{}}
\caption{Container decomposition of the Centinel platform}
\label{tab:containers}\\
\toprule
\textbf{Container} & \textbf{Technology} & \textbf{Problem solved} \\
\midrule
\endfirsthead
\toprule
\textbf{Container} & \textbf{Technology} & \textbf{Problem solved} \\
\midrule
\endhead
\bottomrule
\endlastfoot
Workflow orchestrator & Apache Airflow / Cloud Composer & Schedules discovery, backfills and batch runs; retries failed steps. \\
Scene discovery & Python / STAC client & Finds new Sentinel-1/2 and Landsat scenes over the area of interest. \\
Ingestion service & Python / Cloud Run jobs & Downloads and clips scenes and converts them to cloud-optimized GeoTIFF. \\
Preprocessing service & Python / GDAL on GKE & Cloud masking, calibration, co-registration and tiling; optional SAR module. \\
ML inference service & PyTorch / TorchServe, GPU autoscaling & Runs the multi-label model and produces per-label scores. \\
False-alarm filter & Python rules engine & Applies temporal persistence, spatial coherence and exclusion masks before an alert is issued. \\
Notification service & Python worker & Routes confirmed alerts to subscribers by region, phenomenon and priority. \\
Imagery data lake & Object storage (GCS), COG / Zarr & Single source of raw scenes, analysis-ready tiles and label rasters. \\
Alert database & PostgreSQL + PostGIS & Stores alerts with provenance, field verifications and an append-only audit trail. \\
Alert API & FastAPI / REST & Query, filter and acknowledge alerts; record field verifications. \\
API gateway & Kong / Apigee & Authentication, per-institution roles, rate limiting and access logging. \\
Labeling and ground-truth service & Python / Label Studio & Harmonizes official layers into curated, versioned label sets. \\
Training pipeline & Vertex AI Pipelines / PyTorch & Trains and evaluates joint and independent models with fixed geographic splits. \\
Model registry & MLflow & Versions models with data lineage, metrics and promotion stage. \\
Model and field monitoring & Python / Evidently / Grafana & Tracks drift and field-verified precision; promotes, rolls back or triggers retraining. \\
Web dashboard & React / MapLibre & Map, filters and alert workflow for institutional analysts. \\
Field verification app & React Native, offline-first & Shows nearby alerts and captures geotagged photos and verdicts in areas with poor connectivity. \\
\end{longtable}
}

% ----------------------------------------------------------------
\subsection{Level 3 --- Detection Core}

Figure~\ref{fig:c4-l3} opens the two containers where the central hypothesis is
tested.
Inside the \textbf{ML inference service}, a consumer receives each
\textit{TileReady} event; the tile loader reads the previous-year and current
tiles from the data lake; the feature assembler builds the bitemporal input
tensor (adding SAR channels only if that optional module is enabled); and a
multi-label classifier, with a shared encoder and one head per phenomenon,
produces calibrated scores that the publisher emits as detection events.
The model loader keeps the production model from the registry in memory and
swaps it when a new version is promoted.
Inside the \textbf{false-alarm filter}, each detection passes four sequential
checks: temporal persistence across consecutive observations, spatial coherence
(minimum mapping unit and neighbor agreement), exclusion masks (legal mining
titles, non-forest areas, known crops), and a threshold-and-priority policy that
operators can tune, for example to prioritize parks and indigenous territories.
Only detections that pass all four checks become alerts: the alert writer stores
them with their provenance and emits the confirmation event that drives
notification.

\begin{figure}[H]
\centering
\includegraphics[width=\textwidth]{Cetinel_TDSE-C4_L3_Components.png}
\caption{C4 Level 3: components of the ML inference service and the
false-alarm filter.}
\label{fig:c4-l3}
\end{figure}

% ----------------------------------------------------------------
\subsection{Scalability and Cloud Infrastructure}

The architecture is designed on cloud-native principles.
All compute containers are stateless and deployed via containers
(Docker / Kubernetes), so individual components can be scaled horizontally
without affecting the rest of the pipeline; state lives only in the data lake,
the alert database and the model registry.
Satellite tiles are stored in managed object storage, decoupling data volume
growth from compute costs.
The ML inference service sits behind an auto-scaling endpoint, so peak loads
triggered by new image batches do not require manual capacity planning.
Cross-cutting concerns are handled uniformly: OpenTelemetry for observability,
a secret manager and customer-managed encryption keys, infrastructure as code,
dead-letter queues on every event topic, and data-protection controls to be
reviewed against Colombian Law 1581 of 2012 (Section~\ref{sec:governance}).

Geographic expansion --- from a pilot region to national coverage --- requires
adding new ingestion schedules and storage partitions without architectural
changes.
This separates the \textit{validation problem} (does the approach work?) from
the \textit{scaling problem} (can it cover the whole country?), allowing the
former to be answered cheaply before committing to the latter.


% =============================================================
% PROTOTYPE
% =============================================================

\section{Prototype}
\label{sec:prototype}

The Centinel proof of concept implements and evaluates the satellite-data and machine-learning portion of the proposed architecture. 

Its executed workflow comprises image acquisition through Google Earth Engine, construction of a labeled bitemporal dataset, training of two convolutional classifiers, and evaluation on geographically separated test data. 

A minimal REST interface was implemented and deployed on Google Cloud Run on 17 November 2026. Live HTTP requests exercised the Earth Engine retrieval path and reproduced the Colab outputs for two reference locations. This confirms the behavior of the deployed demonstration container for those requests, not the reliability or operational readiness of the full alert-delivery architecture.

\subsection{Pilot Scope and Research Question}
\label{subsec:prototype-scope}

The original Centinel architecture targets the joint detection of deforestation, illegal alluvial mining, and coca cultivation. A direct comparison between a shared multi-label model and three independent detectors would require spatial reference data for all three phenomena. The necessary layers from IDEAM and UNODC/SIMCI/EVOA were not available within the development period. Therefore, the implemented pilot addresses a narrower and independently testable question: \emph{does adding Sentinel-1 synthetic-aperture radar (SAR) imagery to Sentinel-2 optical imagery improve annual forest-loss detection, particularly when optical observations are partially obscured?}

The pilot is restricted to binary patch-level forest-loss detection. It does not classify mining or coca cultivation, infer the legal status of land-cover change, or generate real-time national alerts. Its purpose is to test a data-fusion design decision and provide reusable technical components for a later multi-phenomenon implementation.

\subsection{Study Area, Satellite Sources, and Dataset Construction}
\label{subsec:prototype-data}

The study region lies within Caquet\'a and Putumayo, Colombia, and is defined in geographic coordinates by the bounding box $(-75.5,\,0.8,\,-74.0,\,1.8)$ (longitude minimum, latitude minimum, longitude maximum, latitude maximum; EPSG:4326). Its approximate surface area is $18{,}000\ \mathrm{km}^{2}$. The region was divided into 24 geographic tiles of $0.25^{\circ}\times 0.25^{\circ}$. Using the Google Earth Engine Python API, the team exported annual composite imagery for 2021, 2022, and 2023, as well as a forest-loss label layer for each tile: 72 image exports and 24 label exports, for 96 completed tasks reported in the repository README.



Each annual image composite contains six bands: four Sentinel-2 surface-reflectance bands (B2, B3, B4, and B8) and two Sentinel-1 radar polarizations (VV and VH). Sentinel-2 observations were screened using the scene-classification layer to exclude clouds, cloud shadows, cirrus, and other unusable pixels; each annual optical composite was formed using a median reducer. Sentinel-1 observations were filtered to interferometric wide-swath mode with both VV and VH polarizations, and were likewise summarized as annual median composites. The working export scale was 20 m per pixel.

The reference labels came from the Hansen/GLAD global forest-change product, using its \texttt{lossyear} band and retaining only areas with at least 30\% tree cover in the year 2000 (Hansen \emph{et al.}, 2013). The label identifies the year of tree-cover loss; it does not identify the cause of the loss or establish that an activity was illegal.

The dataset builder extracted non-overlapping $64\times64$-pixel patches, each approximately 1.28 km on a side. For each location, it paired the previous year's composite with the evaluated year's composite, using the comparisons 2021--2022 and 2022--2023. A patch was labeled positive when at least 2\% of its pixels were marked as forest loss during the evaluated year. Remaining missing pixels were imputed using their respective band means; patches with more than 20\% missing data were discarded. To mitigate class imbalance, the builder retained a random 50\% of negative patches before dataset splitting.

The executed dataset-construction run produced \textbf{11,764 bitemporal patches}, including \textbf{2,461 positive patches (20.9\%)} and 9,303 negative patches. The pipeline saved the image tensors, labels, and geospatial metadata as \texttt{X.npy}, \texttt{Y.npy}, and \texttt{meta.csv}. The positive proportion is specific to the sampled dataset and should not be interpreted as the natural prevalence of forest loss across the region.

\subsection{Model Implementation and Training}
\label{subsec:prototype-model}

The two experimental configurations use the same ResNet-18 backbone, adapted to receive multi-band satellite patches and to output one binary classification logit. The initial convolution was replaced to match the input channel count, the initial max-pooling layer was removed to preserve spatial detail on small patches, and the final layer was replaced with a single-output linear classifier. Both models were trained from scratch rather than initialized with ImageNet weights.

\begin{table}[htbp]
\centering
\small
\caption{Model configurations in the executed comparison.}
\label{tab:prototype-configurations}
\begin{tabular}{|p{0.22\linewidth}|p{0.49\linewidth}|p{0.14\linewidth}|}
\hline
\textbf{Configuration} & \textbf{Input for two consecutive years} & \textbf{Channels} \\ \hline
Optical baseline & Sentinel-2 B2, B3, B4, and B8 in each year & 8 \\ \hline
Optical + SAR fusion & Sentinel-2 B2, B3, B4, and B8 plus Sentinel-1 VV and VH in each year & 12 \\ \hline
\end{tabular}
\end{table}

The experimental code used PyTorch with a class-weighted binary cross-entropy loss (with logits), the AdamW optimizer, a cosine learning-rate schedule, and geometric augmentation through random horizontal and vertical flips and rotations. Training was run for 15 epochs per configuration and repeated with three random seeds on a Google Colab T4 GPU. Per-channel normalization statistics were computed exclusively from the training subset to avoid information leakage.

The train/validation/test design targeted an approximately 70/15/15 division, with group-based splitting assigning complete geographic tiles to each subset. Consequently, the same geographic tile did not occur in both training and held-out subsets, although sample counts varied across seeds. Each seed changed both the model initialization and the geographic partition; the results therefore combine these two sources of variability.

\subsection{Evaluation Protocol and Decision Criteria}
\label{subsec:prototype-evaluation}

For each run, the epoch with the highest validation average precision (AP) was retained, and the binary decision threshold was selected to maximize the F1 score on validation data. The final model was then assessed on the held-out geographic test set using precision, recall, F1, false-positive rate (FPR), AP, and area under the receiver operating characteristic curve (AUROC). Model comparisons used paired differences between the optical and fusion configurations for each seed.

To approximate obscured optical observations, a second evaluation covered a randomly positioned rectangular region in the optical channels of half of the test patches. The occluded pixels were set to the normalized channel mean, while SAR channels were left unchanged. The same occlusion mask was applied to both model configurations. This procedure tests robustness to \emph{synthetic} obstruction; it is not equivalent to evaluation under independently observed cloud patterns.

The predefined criteria for adopting fusion were recall of at least 0.75, precision of at least 0.70, and FPR no greater than 0.20, together with an improvement in mean AP that was positive in at least two of the three paired runs. On clean imagery, fusion yielded an AP of 0.778, versus 0.776 for the optical baseline. Fusion recall was 0.636, below the 0.75 target, while precision was 0.770 and FPR was 0.031. Under synthetic occlusion, AP was 0.661 for fusion and 0.639 for the baseline, a mean difference of $+0.022$. Thus, \textbf{the pilot did not meet its adoption criteria}. The more complete numerical comparison and its interpretation are presented in Section~7.

\subsection{Alert API and Notebook Demonstration}
\label{subsec:prototype-api}

The prototype includes a lightweight FastAPI application (\texttt{04\_api.py}) intended to expose the trained classifier through two HTTP GET endpoints:

\begin{itemize}
    \item \texttt{/health}: returns the service status, selected model configuration, and supported evaluation years.
    \item \texttt{/predict}: accepts \texttt{lat}, \texttt{lon}, and \texttt{year} (2022 or 2023), validates the coordinates against the pilot region, retrieves the previous-year and evaluation-year patches from Earth Engine, normalizes the resulting tensor using the saved checkpoint statistics, and returns a JSON response containing location, compared years, model identifier, forest-loss probability, decision threshold, and binary alert flag.
\end{itemize}

The root path (\texttt{/}) serves a small web page with a map of the pilot area: the user clicks a point inside the pilot box, chooses 2022 or 2023, and receives the forest-loss probability, the decision threshold, and the alert flag, together with a notice that an alert indicates possible forest loss, not illegality, and requires human verification. Two preloaded examples, one with loss (2022) and one without (2023), are provided.

The service keeps the trained fusion model in memory and does not use a database, message queue, or persistent alert store. The \texttt{alert} field is a thresholded model output, not a confirmed incident. The Docker container was built from the repository Dockerfile using Cloud Build and deployed to Google Cloud Run (project \texttt{centinel-tdse}, service \texttt{centinel-api}, region \texttt{us-central1}), with 1~vCPU, 2~GiB RAM, a 300-second request timeout, and port 7860. The image uses Python 3.11 slim and CPU-only PyTorch; the loaded \texttt{fusion} model is a 12-channel ResNet-18 with decision threshold 0.7167.

Hugging Face Spaces was initially considered, but Docker Spaces required a paid plan; Google Cloud Run was selected instead, consistently with the cost model in Section~\ref{sec:business}. The dedicated service account has only Earth Engine Resource Viewer and Service Usage Consumer roles. Its JSON key is stored in Google Secret Manager and injected through the \texttt{EE\_SERVICE\_ACCOUNT\_JSON} environment variable, rather than embedded in source code or the image. The demonstration endpoint permits unauthenticated public requests. Because each prediction retrieves imagery from Earth Engine, no claim is made about measured latency, throughput, availability, concurrency, or scaling.

A companion Colab notebook, \texttt{demo\_centinel\_colab.ipynb}, defines a reproducible demonstration sequence: inspect the tile-export plan; display patch counts and class labels; review stored training logs and comparative metrics; read the pass/fail criteria from \texttt{results.json}; start the API; request \texttt{/health}; query \texttt{/predict} for a labeled positive and a labeled negative patch; and finally request predictions for a balanced demonstration sample of 20 positive and 20 negative test patches. The last step is designed to display a confusion matrix and basic detection counts for the live API interface, rather than to replace the full held-out evaluation.

\textbf{Deployment verification (17 November 2026).} The deployed \texttt{GET /health} endpoint returned \texttt{status=ok}, \texttt{model=fusion}, and supported years 2022 and 2023. A smoke test at latitude 1.3, longitude $-74.75$ (2023) returned probability 0.0185, below the 0.7167 threshold, with no alert. Two reference requests previously used in Colab were then repeated against Cloud Run: the positive location (1.2253, $-74.3909$, 2022) returned 0.9993 and an alert; the negative location (0.8604, $-75.3104$, 2023) returned 0.0338 and no alert. These probabilities matched Colab. The two reference checks and one smoke test demonstrate deployment reproducibility for these cases only; they do not constitute a model evaluation. The notebook's larger demonstration sequence is not claimed as executed in the cloud. No systematic latency, load, concurrency, or scaling measurements were made.

Figure~\ref{fig:frontend} shows the demonstration web interface served at the root path of the deployed service. The dashed rectangle marks the pilot area (Caquet\'a and Putumayo); outside it the model was not trained and the API does not respond. The interface includes two preloaded examples and a notice that the probability is a model output, not the share of forest lost, and that an alert indicates forest loss, not illegality.

\begin{figure}[htbp]
\centering
\includegraphics[width=0.9\linewidth]{Centinel_frontend.png}
\caption{Demonstration web interface of the deployed service (Google Cloud Run). It is a deployment check, not a model evaluation.}
\label{fig:frontend}
\end{figure}


\subsection{Implementation Evidence and Reproducibility}
\label{subsec:prototype-reproducibility}

The implementation is organized into four stages: \texttt{01\_export\_tiles.py} creates the Earth Engine exports; \texttt{02\_build\_dataset.py} assembles the labeled temporal patches; \texttt{03\_train\_eval.py} trains and compares both networks; and \texttt{04\_api.py} defines the alert interface. Shared configuration and image-access utilities are provided by \texttt{config.py} and \texttt{gee\_utils.py}, while \texttt{models.py} defines the ResNet-18 adaptation. The \texttt{make\_synthetic.py} utility is intended for software smoke tests only; synthetic samples are not the basis for the reported detection metrics.

\begin{table}[htbp]
\centering
\small
\caption{Execution evidence supported by the submitted material.}
\label{tab:prototype-verification}
\begin{tabular}{|p{0.28\linewidth}|p{0.60\linewidth}|}
\hline
\textbf{Component} & \textbf{Documented verification status} \\ \hline
Image export & Completed in Earth Engine: 72 Sentinel composites and 24 Hansen/GLAD label exports. \\ \hline
Dataset generation & Executed on satellite-derived data: 11,764 patches, including 2,461 positives. \\ \hline
Training and evaluation & Executed on a Colab T4 GPU: two model configurations across three seeds; metrics reported in the README and paper. \\ \hline
API implementation & FastAPI endpoints \texttt{/health} and \texttt{/predict} exercised over HTTP against the deployed service, including live Earth Engine retrieval. \\ \hline
Cloud deployment & Google Cloud Run (\texttt{us-central1}); Cloud Build from Dockerfile; 1~vCPU, 2~GiB; dedicated service account with two roles and credentials in Secret Manager. Verified 17 November 2026: \texttt{/health} OK; reference positive 0.9993 (alert), negative 0.0338 (no alert), matching Colab. Demonstration URL: \url{https://centinel-api-201209038013.us-central1.run.app}. Service to be deleted after evaluation because the Earth Engine quota is shared. No load or scaling test. \\ \hline
\end{tabular}
\end{table}

The associated scripts and setup instructions are documented in the project repository, \texttt{github.com/A1Daniel1/Centinel}. Numerical results are described as being recorded in \texttt{results/summary.csv} and \texttt{results/results.json}; the submitted README and paper report the key measurements. The checkpoint files are regenerated by the training script and are not included in the repository due to their size. Reproduction requires Earth Engine access and the original exported imagery or repeating the export process.

\subsection{Limitations and Implications for the Full Architecture}
\label{subsec:prototype-limitations}

The prototype supports a limited conclusion about annual forest-loss detection, not about illegal activity. In particular, Hansen/GLAD labels do not identify causes or legal responsibility; yearly composites cannot resolve changes within a year; and artificial rectangular occlusion is only a proxy for tropical cloud conditions. Moreover, the 50\% negative-patch subsampling changed the observed class prevalence, which affects the interpretation of precision; the three random seeds and varying geographic test partitions do not support a statistical-significance claim or extrapolation to the national territory.

The relatively high AUROC reported in Section~7 and the more modest recall suggest that a different threshold may yield a different sensitivity--precision trade-off, but this was not evaluated in the pilot. Equally, the lack of a material clean-image advantage for SAR does not establish that radar is unhelpful in other regions or for shorter temporal windows. The most defensible next step is to validate the original joint-versus-independent hypothesis using suitable official reference data and optical imagery, while retaining SAR as an optional extension. Live Earth Engine requests were successfully exercised through the Cloud Run API on two reference cases and one smoke test. A functioning institutional alert service would still require repeatable operational measurements, access controls, rate limiting, monitoring, and the remaining architecture components.


% =============================================================
% RESULTS
% =============================================================
\section{Results}
\label{sec:results}

\subsection{Dataset and Setup}

The exported area (24 tiles of $0.25^\circ$, 2021--2023) yielded 11{,}764
bitemporal patches of $64\times64$ pixels, of which 2{,}461 (20.9\%) are
positive.
Each configuration was trained for 15 epochs with each of three seeds on a
single GPU.
Test sets contained between 1{,}824 and 1{,}938 patches (148 to 399 positives,
depending on the seed's geographic split).
All figures below come from the pilot's \texttt{results.json}; thresholds were
fixed on the validation split.

\subsection{Quantitative Results}

Table~\ref{tab:results} reports the mean and standard deviation across seeds.
On clean test imagery the two configurations are practically indistinguishable:
average precision is 0.778 for optical plus SAR and 0.776 for optical only,
and F1 is 0.688 and 0.683, respectively.
The per-seed paired difference in average precision is
$+0.003$, $-0.015$, and $+0.017$ (Table~\ref{tab:paired}); its mean
($+0.0016$) is far smaller than the between-seed standard deviation
(0.057--0.067).
Both models rank patches well (AUROC 0.93) at a low false positive rate
(about 0.03).

\begin{table}[ht]
\centering
\caption{Test-set metrics, mean $\pm$ standard deviation over three seeds.
``Occluded'' applies synthetic cloud occlusion to the optical channels of half
of the test patches.}
\label{tab:results}
\small
\begin{tabular}{llcc}
\toprule
Condition & Metric & Optical only & Optical + SAR \\
\midrule
Clean    & Precision & $0.756 \pm 0.094$ & $0.770 \pm 0.129$ \\
         & Recall    & $0.625 \pm 0.022$ & $0.636 \pm 0.044$ \\
         & F1        & $0.683 \pm 0.053$ & $0.688 \pm 0.044$ \\
         & FPR       & $0.030 \pm 0.007$ & $0.031 \pm 0.016$ \\
         & AP        & $0.776 \pm 0.067$ & $0.778 \pm 0.057$ \\
         & AUROC     & $0.931 \pm 0.008$ & $0.931 \pm 0.011$ \\
\midrule
Occluded & Precision & $0.769 \pm 0.097$ & $0.785 \pm 0.111$ \\
         & Recall    & $0.429 \pm 0.063$ & $0.454 \pm 0.068$ \\
         & F1        & $0.550 \pm 0.073$ & $0.567 \pm 0.053$ \\
         & FPR       & $0.019 \pm 0.005$ & $0.021 \pm 0.011$ \\
         & AP        & $0.639 \pm 0.103$ & $0.661 \pm 0.093$ \\
         & AUROC     & $0.846 \pm 0.026$ & $0.860 \pm 0.034$ \\
\bottomrule
\end{tabular}
\end{table}

\begin{table}[ht]
\centering
\caption{Paired per-seed difference (optical + SAR minus optical only).}
\label{tab:paired}
\small
\begin{tabular}{lccc}
\toprule
 & Seed 0 & Seed 1 & Seed 2 \\
\midrule
$\Delta$AP, clean    & $+0.003$ & $-0.015$ & $+0.017$ \\
$\Delta$F1, clean    & $-0.020$ & $+0.017$ & $+0.019$ \\
$\Delta$AP, occluded & $+0.022$ & $+0.009$ & $+0.036$ \\
\bottomrule
\end{tabular}
\end{table}

Under synthetic occlusion, both models degrade: average precision falls from
0.776 to 0.639 (optical only) and from 0.778 to 0.661 (optical plus SAR), and
recall drops to 0.43 and 0.45.
The fusion model is ahead in all three seeds ($\Delta$AP of $+0.022$ on
average), a consistent but small advantage that is of the same order as the
between-seed variation.

\subsection{Evaluation Against the Decision Criteria}

Against the criteria of Section~\ref{sec:gap}, the optical-plus-SAR
configuration meets the precision requirement (0.770 $\geq$ 0.70) and the false
positive rate requirement (0.031 $\leq$ 0.20), but \textbf{not} the recall
requirement (0.636 $<$ 0.75).
The criterion that fusion improve average precision is satisfied only in the
formal sense (positive mean and positive in two of three seeds); the margin is
negligible and cannot be distinguished from seed-to-seed variation.
Fusion is also more robust to occlusion in all three seeds, with a small effect.
The pilot therefore \emph{does not meet} its go criteria: the evidence supports
neither adopting nor discarding SAR fusion, and the limiting factor on this
dataset is recall rather than the choice of input.

\subsection{Implications for the Original Hypothesis}

Because the adjusted question was not answered in favor of fusion, the
pilot gives no basis for making SAR ingestion a requirement of the first
validation of the joint-detection hypothesis.
Radar adds data volume, co-registration effort, and processing cost, and on
annual composites it produced no meaningful gain on clean imagery and only a
small gain under occlusion.
The pilot also does not weaken the original hypothesis, since it did not
test it.
Since radar adds cost and no benefit was demonstrated, we decide, as a matter
of cost under uncertainty and not as a result of the pilot, to test joint
versus independent detection with optical imagery only, keeping SAR as an
optional module to be revisited if cloud cover proves limiting at larger scale
or with denser time series.
The pilot does not show that radar is useless: it shows that, on annual
composites and this dataset, its benefit could not be demonstrated.

\subsection{Interpretation and Limits}

Three observations bound the interpretation.
First, the annual cloud-masked median composites already remove most cloud
effects before the model sees the data, which may leave little for radar to
add on clean imagery; this is a hypothesis, not something the pilot tested.
Earlier work on optical--SAR fusion in the Brazilian Amazon reports gains from
fusion~\cite{ferrari2023}; it fuses bitemporal image pairs, whereas the pilot
used annual median composites, and whether that difference explains the
contrast was not tested.
Second, the high AUROC alongside moderate recall indicates that the model ranks
patches well but that the operating point selected for F1 on the validation
split trades recall for precision; a threshold tuned for recall, or the
label noise inherent to a global 30\,m forest-loss product, could shift this,
but neither was explored.
Third, with three seeds, geographic splits that change with the seed, and no
significance test, none of the differences above should be read as established.
The pilot evaluates the detection model separately from the deployment check. The API was deployed on Google Cloud Run and two reference queries reproduced the Colab probabilities (0.9993 with alert and 0.0338 without alert); a third request served as a smoke test of the Earth Engine connection. This is not a new evaluation of the model. Latency, throughput, concurrency, load, and scaling were not systematically measured.


% =============================================================
% CONCLUSIONS AND FUTURE WORK
% =============================================================
\section{Business Case, Governance and Responsibility}
\label{sec:business}

This section addresses the aspects of the proposal that the pilot does not
evaluate empirically: what the system would cost, how its value would be
measured, how it would be governed and secured, and what responsibilities it
raises, and it ends with a phased roadmap. Each figure is labeled as
\emph{measured} in the pilot, \emph{derived} by arithmetic from a measurement, or
an \emph{assumption}. Prices were read from the providers' public pages in
November 2026 and change over time.

\subsection{Cost Model and Value Metrics}
\label{sec:costs}

Satellite data from Copernicus and USGS/NASA is free of charge, so the cost of
the platform is driven by processing, storage, human labeling, and the team
that operates it. We model the annual cost as
\begin{equation}
C = C_{\text{fixed}} + c_{\text{proc}}\,A\,f + c_{\text{stor}}\,V + c_{\text{lab}}\,N_{\text{lab}},
\end{equation}
where $A$ is the monitored area, $f$ the revisit frequency, $V$ the stored
volume, and $N_{\text{lab}}$ the number of labeled samples per year.
Table~\ref{tab:costs} lists the drivers, how each scales, and what we know about
it.

\begin{table}[H]
\centering
\caption{Cost drivers of the platform (parametric model; the Earth Engine row and
the other priced items are detailed in Section~\ref{sec:cost-estimate}).}
\label{tab:costs}
\begin{tabularx}{\linewidth}{l l X}
\toprule
\textbf{Driver} & \textbf{Scales with} & \textbf{Comment} \\
\midrule
Imagery acquisition & -- & Source imagery is free (Copernicus, USGS/NASA); processing through Earth Engine may carry a license cost. \\
Processing (tiling, inference) & $A \cdot f$ & Stateless containers; cost grows linearly with area and revisit, and can be reduced by restricting inference to cloud-free tiles. \\
Storage (scenes, tiles, detections) & $V$ & Dominated by raw scenes; mitigated by keeping only derived tiles after a retention period. \\
Labeling and field verification & $N_{\text{lab}}$ & Expected to be the main recurring cost, since expert labels are scarce and were the limiting factor in this work. \\
Platform operation & fixed & Gateway, identity provider, monitoring, and a small operations team. \\
\bottomrule
\end{tabularx}
\end{table}

\subsubsection{Cost Estimate}
\label{sec:cost-estimate}

\paragraph{Earth Engine.}
The pilot measured about 352 Earth Engine compute-seconds (EECU-s) to export one
annual six-band composite of one $0.25^\circ$ tile (roughly 770\,km$^2$).
Colombia's continental territory is about 1.14 million\,km$^2$, or about 1{,}500
such tiles, so a national pass costs about $1{,}500\times352 \approx 528{,}000$
EECU-s, or \textbf{about 147 EECU-hours} (derived from a single measurement).
Table~\ref{tab:ee} translates this into yearly volumes and cost under the two
licensing situations, using the public pricing~\cite{ee_pricing} and
noncommercial quotas~\cite{ee_noncommercial}. The calculation assumes the cost
per composite does not fall with shorter time windows; shorter windows use fewer
images, so this is likely conservative, but it was not tested.

\begin{table}[H]
\centering
\caption{Earth Engine volume and cost by processing cadence for national
coverage (derived; one national pass $\approx$ 147 EECU-hours).}
\label{tab:ee}
\small
\begin{tabular}{lrp{4.1cm}p{4.3cm}}
\toprule
Cadence & EECU-h/year & Noncommercial (Community tier, 150 EECU-h/month) & Commercial (Basic plan, \$500/month incl.\ 100 batch EECU-h/month) \\
\midrule
Annual    & 147   & Fits (one pass)                  & about \$6{,}000/year (within included hours) \\
Quarterly & 588   & Fits (147 per quarter)           & about \$6{,}000/year (within included hours) \\
Monthly   & 1{,}764 & About 147 per month: barely fits the 150 limit & about \$6{,}226/year (\$6{,}000 plus 564 extra EECU-h at \$0.40) \\
\bottomrule
\end{tabular}
\end{table}

Whether an operational service run by or for a government agency qualifies as
noncommercial is not stated in the quota documentation we consulted, and we do
not assume that it does. Under a commercial license the monthly platform fee
dominates; usage is a small share.

\paragraph{Other items.}
Table~\ref{tab:cost} lists the remaining items.

\begin{table}[H]
\centering
\caption{Other cost items. ``Not estimated'' means we did not obtain a
defensible figure.}
\label{tab:cost}
\small
\begin{tabular}{p{2.6cm}p{5.4cm}p{3.7cm}p{1.9cm}}
\toprule
Item & Basis & Estimate & Status \\
\midrule
Model training (GPU) &
Pilot: six runs took under ten minutes on one T4. National data would be about
62$\times$ larger (1{,}500 vs.\ 24 tiles), assuming linear scaling. One
g4dn.xlarge costs \$0.526/hour on demand and about \$0.21--0.24 as spot~\cite{g4dn_price} &
Up to about 10 hours, or about \$5.50 on demand, per equivalent set of six runs;
about \$110 for 20 such sets & Assumption \\
Composite storage &
6 bands $\times$ 4 bytes at 20\,m is about 60\,KB/km$^2$, or about 68\,GB per national
annual composite (uncompressed). Earth Engine storage list price is about
\$0.026/GB-month, used as a proxy for object storage &
About \$1.80/month per year of data retained & Derived \\
Alert API serving &
Assumed 8\,s of one vCPU and 0.5\,GiB per request; Cloud Run list prices of
\$0.000024/vCPU-s and \$0.0000025/GiB-s~\cite{cloudrun_pricing} &
About \$0.0002 per request; roughly 22{,}000 requests/month fall in the free tier &
Assumption (latency not measured) \\
Database, orchestrator, observability, security tooling &
Not priced. Managed options vary widely in cost &
Not estimated & --- \\
Staff, operations, legal review &
Not priced; likely the largest item &
Not estimated & --- \\
\bottomrule
\end{tabular}
\end{table}

\paragraph{Reading of the numbers.}
Compute and storage for a national annual pass are small. The decisions that
move the budget are the Earth Engine license status, the managed services and
staff we did not price, and whether the cadence needs to go below annual.

Because the pilot was scoped to deforestation and the architecture is cloud
native and stateless, the pilot cost is small and bounded, whereas the national
cost depends mainly on the license status, on $N_{\text{lab}}$ and on the operations team. We do not
estimate savings: the value of early detection depends on interventions that the
platform does not control.
Instead, we propose operational metrics that the institution can measure from
the first deployment:
\begin{itemize}
    \item \textbf{Detection-to-alert latency}: time from scene availability to an alert in the Alert API.
    \item \textbf{Verified-alert precision}: share of alerts that field verification confirms, recorded through the field application.
    \item \textbf{Recall against an independent reference}, for example official deforestation maps, at a fixed area.
    \item \textbf{Hectares monitored per verification-officer-day}, as a measure of how much the system reduces manual review.
    \item \textbf{Cost per verified alert}, combining processing, labeling and verification effort.
\end{itemize}
These metrics tie the technical evaluation of Section~\ref{sec:results} to the
decision an institution would face, namely whether the cost per verified alert
is acceptable.

\subsection{Governance and Security}
\label{sec:governance}

\textbf{Security status of the deployed demonstration.} The Cloud Run service is publicly accessible without authentication (\texttt{--allow-unauthenticated}). Its dedicated service account is limited to Earth Engine Resource Viewer and Service Usage Consumer; the service-account JSON key is held in Secret Manager and injected as \texttt{EE\_SERVICE\_ACCOUNT\_JSON}. Authentication, rate limiting, and application-level monitoring/observability remain proposed controls, not implemented features. The public service is intended only for the demonstration and will be deleted after evaluation because the Earth Engine quota is shared.


\textbf{Governance as error correction.}
Following the view that institutions should be designed so that mistaken
decisions can be detected and replaced~\cite{popper1945}, we treat the
architecture and the detection model as revisable proposals, within the decision
rights and accountability that an IT governance framework
requires~\cite{weill2004}. The policy for Centinel specifies the following.
\begin{itemize}
    \item \textbf{Hypothesis and refuting evidence.} The hypothesis is that one joint multi-label model detects the three phenomena better than three independent detectors. Evidence against it would be a joint model that does not match independent detectors in recall at equal precision on the same labeled region.
    \item \textbf{Evaluation as organized criticism.} The decision criteria of Section~\ref{sec:pilot-criteria} were fixed before the experiment, and the pilot's result, that the criteria were not met, is reported as such. The scope was then revised instead of defending the original plan: since label layers for mining and coca could not be obtained, the pilot was reduced to deforestation. We regard this as the process working as intended.
    \item \textbf{Authority and prohibited actions.} The system may generate and rank alerts; it may not trigger operations, sanctions or communications to third parties, and no alert is treated as proof of illegality.
    \item \textbf{Roles.} A \emph{data owner} (an environmental authority) decides which layers may be used for training; a \emph{model owner} approves each model version before the registry promotes it; a \emph{verification lead} accepts or rejects field confirmations. In a real deployment, an \emph{independent reviewer} outside the team that builds the model, for example an environmental authority or a technical committee, would examine disputed decisions; this role was not applied in the pilot.
    \item \textbf{Stop conditions, appeal and revision.} The model owner or the reviewer suspends the system if verified-alert precision, measured over the last $N$ field verifications, falls below a threshold $P$, or if the reference data prove unreliable. The threshold $P$ and the window $N$ are set with the data owner after an initial field-verification phase; before that, no value can be justified from pilot data. Contested alerts feed the evaluation set. The policy has an owner and a scheduled review date, and each revision records its reasons.
\end{itemize}
Retraining is triggered by monitoring but promoted only after the model owner
reviews the evaluation on a held-out region, so that no model reaches
production without human approval.

\textbf{Governance and management.} In data-management terms,
governance sets decision rights and accountability, and management executes
within them~\cite{dama2017}. In Centinel the data owner and the model owner
govern, in the sense that they decide which layers and which model versions are
admissible, whereas the verification lead and the operations team manage, in
the sense that they run the pipeline and confirm alerts.

\textbf{Quality cycle.} The pilot followed a plan--do--check--act cycle, which
is the quality cycle that data-quality guidance builds on~\cite{dama2017}:
\emph{plan}, the decision criteria were fixed before the experiment;
\emph{do}, the pilot was run on open data; \emph{check}, the results were
compared with the criteria (recall 0.636 against a target of 0.75); and
\emph{act}, the scope was reduced and the roles and stop condition above were
defined. A production deployment would repeat this cycle on a schedule, not only
once.

\textbf{Governance of reference data.} The main obstacle of this work was not the
model but access to ground truth: the IDEAM, UNODC-SIMCI and EVOA layers were
not obtained. A deployable policy therefore needs data-sharing agreements with
those institutions, a named custodian of the labels, and a rule for how they are
updated; until then, the platform can only be trained on layers it is allowed to
use.

\textbf{Data quality.} The quality of the reference labels and of the imagery
bounds what any model can achieve, so we select the dimensions of the ISO/IEC
25012 data-quality model~\cite{iso25012} that matter for Centinel and attach an
indicator and a responsible role to each (Table~\ref{tab:dq}). The indicators
are proposals for a production phase; the pilot did not measure them.

\begin{table}[H]
\centering
\caption{Data-quality dimensions selected for Centinel (proposed indicators).}
\label{tab:dq}
\small
\begin{tabularx}{\linewidth}{l X X l}
\toprule
\textbf{Dimension} & \textbf{Where it matters} & \textbf{Indicator} & \textbf{Role} \\
\midrule
Accuracy & Forest-loss labels from global products omit small-scale clearing & Agreement with official layers on a held-out region & Data owner \\
Completeness & Cloud cover leaves pixels without data & Share of valid pixels per tile and period & Model owner \\
Currentness & Labels are annual, alerts are meant to be earlier & Delay between an event and the availability of its reference label & Data owner \\
Consistency & Optical and radar inputs must share grid and projection & Share of tiles passing alignment checks & Model owner \\
Traceability & An alert must be explainable after the fact & Share of alerts linked to a scene, a model version, and a label version & Model owner \\
Confidentiality & Alert coordinates are sensitive & Denied and logged access attempts to exact coordinates & Data owner \\
\bottomrule
\end{tabularx}
\end{table}

As an example of how each indicator would be specified, the stop condition
of this section uses \emph{verified-alert precision}: its quality attribute is
accuracy; its variables are the number of alerts confirmed in the field,
$C$, and the number of alerts verified, $V$, over the last $N$ verifications;
its formula is $\hat{p} = C/V$; its target is the threshold $P$, set with the
data owner after an initial verification phase; and a value below $P$ is read as
a reason to suspend the system, not as proof of a model failure, because the
reference data may also be at fault.

\textbf{Security design.} Following secure-design
principles~\cite{saltzer2009}, every request to the Alert API passes through the
API gateway and the identity provider (complete mediation), access is denied
unless a role grants it (fail-safe defaults), and alert coordinates are visible
only to authorized officers. Alert state changes and field confirmations are
logged for auditing.

\textbf{Proposed production baseline.} The pilot does not implement the
following controls; they are the baseline we propose for a production phase.
\begin{itemize}
    \item \textbf{Authentication and transport:} API keys or OAuth~2 for
    institutional clients, TLS on every endpoint, and per-client rate limits,
    which also protect the Earth Engine quota.
    \item \textbf{Secrets and access:} cloud secret storage for service-account
    keys, least-privilege roles, and separate credentials per environment.
    \item \textbf{Operational metrics:} request latency, error rate, Earth
    Engine quota consumption, and fraction of requests rejected for missing data.
    \item \textbf{Model monitoring:} alert rate per region and period, input
    missing-data fraction, and probability distribution drift, as early signs
    that the model no longer matches the data. Labeled performance can only be
    re-measured when new reference data arrive.
    \item \textbf{Audit log:} the log of alert state changes is extended to
    record who queried what, to support later review of how an alert was
    produced.
    \item \textbf{Service objectives:} availability and latency targets should be
    set with the institutional users; we do not propose figures that the pilot
    did not measure.
\end{itemize}

\textbf{Publication policy and legal review.} The institution should define who
receives exact coordinates, who receives aggregated views, and for how long
alerts are retained. Alerts inside indigenous territories and protected areas
may require prior consultation and community agreements, and location data
linked to identifiable people falls under Colombia's personal-data law (Law 1581
of 2012). Both points need legal review; we do not assert how they apply.

Table~\ref{tab:risks} summarizes the main risks.

\begin{table}[H]
\centering
\caption{Risk matrix (qualitative; likelihood and impact are our judgment, not measured).}
\label{tab:risks}
\begin{tabularx}{\linewidth}{X c c X}
\toprule
\textbf{Risk} & \textbf{Likelihood} & \textbf{Impact} & \textbf{Mitigation} \\
\midrule
Label scarcity prevents training for mining and coca (\emph{occurred}) & High & High & Request IDEAM, UNODC-SIMCI and EVOA layers; start with deforestation only. \\
Recall below target (\emph{occurred}: 0.64 vs.\ 0.75) & High & High & Improve labels and architecture before any deployment; report recall with every alert statistic. \\
Labels reflect a definition of forest loss, not illegality & High & High & Human confirmation; treat alerts as leads, not findings; validate against official layers. \\
Poor generalization to other regions or years & Medium & High & Evaluate on held-out regions before expanding; regional deployment first. \\
False positives lead to unwarranted field operations & Medium & High & Confidence threshold, human confirmation before any action. \\
Missed detections under persistent cloud cover & High & Medium & Optional radar module; cloud-aware scheduling. \\
Model drift as land cover changes & Medium & Medium & Monitoring and approved retraining. \\
Leak of alert coordinates & Low & High & Proposed role-based access, audit log, encryption at rest and in transit; demo endpoint currently public and unauthenticated. \\
Dependence on free data providers and licensing (Earth Engine terms, quotas) & Medium & Medium & Multiple sources (Sentinel-2, Landsat); confirm license status before operation. \\
\bottomrule
\end{tabularx}
\end{table}

\subsection{Ethical, Organizational and Philosophical Implications}
\label{sec:ethics}

\textbf{An alert is not evidence.} The system produces a probabilistic
statement about pixels, not a finding of illegality: a clearing may be
authorized, and a mine or a crop may belong to an activity that is lawful or
under a substitution agreement. Treating a detection as proof would shift
responsibility from officials to an algorithm. For this reason the architecture
keeps a human confirmation step, and the explainability required for judicial
use remains future work.

\textbf{Who bears the cost of errors.} False positives fall on people: rural and
campesino communities, and Indigenous and Afro-descendant territories, who may
face operations triggered by a mistaken alert. The pilot's recall of 0.636 shows
that the model would also miss a large part of real events. Both error types
should be reported openly to the institutions that use the system.

\textbf{Dual use.} A tool built to protect ecosystems can also support
surveillance of communities. Access control, purpose limitation, and a policy on
who may request which alerts are therefore part of the design, not an
afterthought, in line with the governance function of the NIST AI risk
management framework~\cite{nistairmf2023}.

\textbf{Bias in the reference data.} Labels derived from global forest-loss
products encode a definition of forest and a detection method; a model trained on
them inherits their omissions, for example small-scale clearing under cloud.
Claims about performance should be tied to the reference used.

\textbf{Organizational implications.} Adopting the platform requires
inter-institutional agreements on data sharing, clear responsibility among
environmental authorities, the prosecutor's office and security forces, and
training for the officers who verify alerts. Without these, a technically sound
system would produce alerts that no one is accountable for acting on.
\subsection{Roadmap}
\label{sec:roadmap}

Timings are indicative planning assumptions, not derived from the pilot.

\begin{table}[H]
\centering
\caption{Phased roadmap with decision gates.}
\label{tab:roadmap}
\small
\begin{tabular}{p{1.8cm}p{1.9cm}p{4.2cm}p{4.0cm}}
\toprule
Phase & Timing & Objective & Exit gate \\
\midrule
0. Pilot & Completed (2026) & Forest-loss pipeline on open data; SAR fusion question & Not met: recall 0.64 vs.\ 0.75; hypothesis untested \\
1. Test the hypothesis & Months 0--3 after approval & Obtain ground truth; compare joint vs.\ independent detection, optical only & Joint model lowers false alarms in co-occurrence zones, paired across seeds on a fixed geographic split \\
2. Production MVP & Months 3--9 & Ingestion, registry, alert database, orchestration, security and observability baseline; regional cloud deployment & Service objectives agreed and met in one region \\
3. Scale and integrate & Months 9--18 & Additional regions or national coverage; integration with institutional systems; governance in operation & Institutional adoption decision \\
\bottomrule
\end{tabular}
\end{table}

\section{Conclusions and Future Work}
\label{sec:conclusions}

\subsection{Contributions}

This paper proposed Centinel, a cloud reference architecture for the early
detection of illegal territorial activities from satellite imagery in Colombia,
organized around one central hypothesis: that jointly detecting deforestation,
illegal mining, and coca cultivation lowers false alarms relative to three
independent detectors.
The ground truth needed to test it (IDEAM, UNODC/SIMCI, EVOA) could not be
obtained within the project timeline, so the question and the minimum viable
product were adjusted to open data.
The resulting pilot tested whether adding Sentinel-1 radar to Sentinel-2 imagery
improves annual forest-loss detection and robustness to cloud occlusion.
On 11{,}764 patches over a $1.5^\circ\times1^\circ$ area of Caquetá and
Putumayo, radar fusion gave no meaningful gain on clean imagery, a small
consistent gain under synthetic occlusion, and recall of 0.64 against a target
of 0.75.
The adjusted question was therefore not answered in favor of fusion.
This does not close the project: the original hypothesis is still viable and
untested, and, since radar showed no
demonstrated benefit, we propose to test it first with optical imagery only.
The pilot also delivers a working, reproducible pipeline from imagery to
detection that the next validation can reuse.

\subsection{Final Context}

The results concern forest-loss detection measured against a global
tree-cover-loss product; they do not show that illegal activity can be
detected or attributed.
Any institutional use, for example by environmental authorities, would require
validation against official deforestation data and an assessment of cause and
legality, which are outside the scope of this pilot.

\subsection{Future Work}

The priority is the original premise, tested with the pipeline built for the
pilot but without radar:

\begin{itemize}
    \item \textbf{Test joint versus independent detection (priority).}
    Obtain spatial ground truth from IDEAM (deforestation), UNODC/SIMCI (coca
    cultivation), and the competent mining authority (alluvial mining), and
    compare a shared multi-label model with three independent detectors on
    Sentinel-2 imagery only, in a region where the phenomena co-occur.
    Planned metrics are those of the pilot (precision, recall, F1, false
    positive rate, average precision), with false alarm rate in co-occurrence
    zones as the primary outcome, repeated over seeds and with a fixed
    geographic split to allow a paired significance test.
    \item \textbf{Keep SAR optional.} Revisit radar if cloud cover limits
    optical detection at larger scale or with intra-annual time series.
    \item \textbf{Close the recall gap:} tune the decision threshold for
    recall, evaluate pretrained and transformer-based backbones, and test
    higher-resolution or denser temporal inputs.
    \item \textbf{Official labels:} replace Hansen / GLAD labels with the
    official IDEAM deforestation product.
    \item \textbf{Harden the deployed demonstration:} replace public unauthenticated access with authentication and rate limiting, add application observability, and test latency, load, concurrency and regional reliability before operational use.
    \item \textbf{Toward full deployment:} intra-annual time series for
    near-real-time alerts, the remaining containers of the architecture,
    explainability suitable for judicial evidence, and governance and data
    ownership arrangements for indigenous territories.
\end{itemize}


% =============================================================
% BIBLIOGRAPHY
% =============================================================
\section*{Declaration on the Use of Artificial Intelligence}
The authors used artificial intelligence tools, specifically ChatGPT and Claude, as support for writing and reviewing this document. The texts were first drafted from the authors' own ideas and judgment; the tools were then used to identify aspects to improve, mainly grammar correction, clarity, organization and the way ideas are expressed, and to help verify references and update public figures. The suggestions were not incorporated automatically: the authors reviewed each one, checked figures and sources against the original documents, and accepted or rejected the changes according to their own criteria.

The technical development of the work, including the architecture design, the pipeline implementation, the experiments, the results, the analysis and the methodological decisions, was carried out by the authors.

\begin{thebibliography}{99}

\bibitem{ideam_smbyc}
IDEAM. \textit{Alertas Tempranas de Deforestación}. Sistema de Monitoreo de
Bosques y Carbono (SMByC), 2016.
\url{http://www.ideam.gov.co/web/ecosistemas/alertas-tempranas-por-deforestacion}

\bibitem{eltiempo_deforestacion2025}
El Tiempo. \textit{La deforestación en Colombia aumentó 5,2\% en 2025}. 2025.
\url{https://www.eltiempo.com/vida/medio-ambiente/la-deforestacion-en-colombia-aumento-5-2-en-2025-pero-el-gobierno-celebra-cque-cifras-estan-por-debajo-de-lo-pactado-en-el-plan-de-desarrollo-3575461}

\bibitem{simci2024}
Forbes Colombia. \textit{Cultivos de coca llegaron a 261.000 hect\'areas en 2024
tras aumentar 3,5\%}. June 25, 2026 (report \textit{Monitoreo de territorios con
presencia de cultivos de coca 2024}, UNODC and Ministerio de Justicia).
\url{https://forbes.co/2026/06/25/actualidad/cultivos-de-coca-llegaron-a-261-000-hectareas-en-2024-tras-aumentar-35/}

\bibitem{siima2025}
ColombiaCheck. \textit{Honorio Henr\'iquez habla de 330.000 hect\'areas de coca
cultivadas en Colombia, pero cifras oficiales son menores}. 2026 (reports SIIMA,
Polic\'ia Nacional, and El Tiempo, Jan.\ 2026: 262,179\,ha, cut-off August 2025).
\url{https://colombiacheck.com/chequeos/honorio-henriquez-habla-de-330000-hectareas-de-coca-cultivadas-en-colombia-pero-cifras}

\bibitem{minenergia_evoa2025}
Ministerio de Minas y Energ\'ia y UNODC. \textit{Colombia registr\'o 118.410
hect\'areas con explotaci\'on de oro de aluvi\'on en 2025}. Informe EVOA 2025,
October 2, 2026.
\url{https://minenergia.gov.co/es/sala-de-prensa/noticias-index/colombia-registro-118410-hectareas-con-explotacion-de-oro-de-aluvion-en-2025-segun-el-informe-anual-del-gobierno-nacional-y-unodc/}

\bibitem{minjusticia_coca2023}
Ministerio de Justicia y del Derecho. \textit{Resultados del monitoreo a
cultivos de coca hecho en 2023}. UNODC/SIMCI, 2023.
\url{https://www.minjusticia.gov.co/Sala-de-prensa/Paginas/Resultados-monitoreo-cultivos-2023-no-son-ajenos-a-presion-global-por-aumento-de-demanda.aspx}

\bibitem{hansen2013}
M.~C. Hansen, P.~V. Potapov, R.~Moore, M.~Hancher, S.~A. Turubanova,
A.~Tyukavina, et al. \textit{High-Resolution Global Maps of 21st-Century Forest
Cover Change}. Science, 342(6160):850--853, 2013.
\url{https://doi.org/10.1126/science.1244693}

\bibitem{mongabay_siona2023}
Mongabay Latam. \textit{Deforestación, minería ilegal y narcotráfico acorralan
a resguardos indígenas de la Amazonía de Colombia}. 2023.
\url{https://es.mongabay.com/2023/11/deforestacion-mineria-ilegal-narcotrafico-acorralan-resguardos-indigenas-amazonia-colombia/}

\bibitem{mongabay_riopure2026}
Mongabay Latam. \textit{Minería ilegal en Colombia: deforestación por el oro
avanza en el Parque Río Puré y pone en riesgo a pueblos indígenas en aislamiento}.
2026.
\url{https://es.mongabay.com/2026/07/mineria-ilegal-colombia-deforestacion-oro-rio-pure/}

\bibitem{riocampo2023}
S.~L. Riocampo, J.~A. Correa García, A.~Gómez Luévano, and G.~S. Bermúdez
García. \textit{Inteligencia artificial para la detección de la minería ilegal,
deforestación y cultivos ilícitos por medio de imágenes satelitales}.
Encuentro Internacional de Educación en Ingeniería, ACOFI, 2023.
\url{https://acofipapers.org/index.php/eiei/article/view/2973/2315}

\bibitem{comimo_urosario}
Universidad del Rosario. \textit{CoMiMo: Detección de minería ilegal mediante
imágenes satelitales}. Based on Saavedra-Romero (2017), 2017.
\url{https://pure.urosario.edu.co/es/projects/miner\%C3\%ADa-ilegal/}

\bibitem{smbyc_atd_github}
SMByC. \textit{ATD: Alertas Tempranas de Deforestación}. Source code repository,
IDEAM, 2016.
\url{https://github.com/SMByC/ATD}

\bibitem{facsat1}
Fuerza Aérea Colombiana. \textit{FACSAT-1: Detección de patrones de minería
ilegal a cielo abierto utilizando redes neuronales}. 2023.
\url{https://poderespacial.fac.mil.co/es/facsat/deteccion-de-patrones-de-mineria-ilegal-cielo-abierto-utilizando-redes-neuronales}

\bibitem{unodc_evoa2022}
UNODC. \textit{Explotaci\'on de oro de aluvi\'on: evidencias a partir de
percepci\'on remota 2022} (resumen ejecutivo). Oficina de las Naciones Unidas
contra la Droga y el Delito, Colombia, 2023.
\url{https://www.unodc.org/documents/colombia/2023/noviembre-11/Resumen_Ejecutivo_EVOA_2022.pdf}

\bibitem{reiche2021radd}
J.~Reiche, A.~Mullissa, B.~Slagter, et al. \textit{Forest disturbance alerts
for the Congo Basin using Sentinel-1}. Environmental Research Letters,
024005, 2021.
\url{https://doi.org/10.1088/1748-9326/abd0a8}

\bibitem{gfw_radd}
Global Forest Watch. \textit{New Radar Alerts Monitor Forests Through the
Clouds}. World Resources Institute, 2021.
\url{https://www.globalforestwatch.org/blog/data/radd-radar-alerts}

\bibitem{gfw_integrated}
Global Forest Watch. \textit{Looking for the Quickest Signal of Deforestation
or Vegetation Disturbance Globally? Turn to GFW's Integrated Disturbance
Alerts}. 2026.
\url{https://www.globalforestwatch.org/blog/data/integrated-deforestation-alerts}

\bibitem{ferrari2023}
F.~Ferrari, M.~P. Ferreira, and R.~Q. Feitosa. \textit{Fusing Sentinel-1 and
Sentinel-2 images with transformer-based network for deforestation detection
in the Brazilian Amazon under diverse cloud conditions}. ISPRS Annals of the
Photogrammetry, Remote Sensing and Spatial Information Sciences,
X-1/W1-2023, 2023.
\url{https://isprs-annals.copernicus.org/articles/X-1-W1-2023/999/2023/}


\bibitem{popper1945}
K. Popper. \textit{The Open Society and Its Enemies}. Routledge, London, 1945.

\bibitem{weill2004}
P. Weill and J. Ross. \textit{IT Governance: How Top Performers Manage IT
Decision Rights for Superior Results}. Harvard Business School Press, 2004.

\bibitem{saltzer2009}
J. H. Saltzer and M. F. Kaashoek. \textit{Principles of Computer System Design:
An Introduction}. Morgan Kaufmann, 2009.

\bibitem{nistairmf2023}
National Institute of Standards and Technology. \textit{Artificial Intelligence
Risk Management Framework (AI RMF 1.0)}. NIST AI 100-1, 2023.

\bibitem{ee_pricing}
Google Cloud. \textit{Google Earth Engine pricing}. Accessed November 2026.
\url{https://cloud.google.com/earth-engine/pricing}

\bibitem{ee_noncommercial}
Google. \textit{Earth Engine noncommercial quota tiers}. Accessed November 2026.
\url{https://developers.google.com/earth-engine/guides/noncommercial_tiers}

\bibitem{cloudrun_pricing}
Google Cloud. \textit{Cloud Run pricing}. Accessed November 2026.
\url{https://cloud.google.com/run/pricing}

\bibitem{g4dn_price}
DoiT. \textit{AWS EC2 g4dn.xlarge pricing in us-east-1} (third-party price
aggregator). Accessed November 2026.
\url{https://www.doit.com/compute/spot/us-east-1/g4dn.xlarge}

\bibitem{minambiente_13}
Ministerio de Ambiente y Desarrollo Sostenible. \textit{El 13\,\% de la deforestaci\'on anual del pa\'is est\'a vinculada a cultivos de uso il\'icito}. Press release, September 8, 2023.
\url{https://www.minambiente.gov.co/el-13-de-la-deforestacion-anual-del-pais-esta-vinculada-a-cultivos-de-uso-ilicito/}

\bibitem{dama2017}
DAMA International. \textit{DAMA-DMBOK: Data Management Body of Knowledge}, 2nd ed.
Technics Publications, 2017.

\bibitem{iso25012}
ISO/IEC 25012:2008. \textit{Software engineering --- Software product Quality Requirements and Evaluation (SQuaRE) --- Data quality model}. International Organization for Standardization, 2008.

\end{thebibliography}

\end{document}